% Full predictable-drift, quadratic-variation, adaptive-count and stopping proofs.

\begin{annexitem}{annex:p15.doob-decomposition}{thm:p15-doob-decomposition}{Every drift term and uniqueness}
\begin{proof}
The increment \(Y_n-Y_{n-1}\) is integrable, so its conditional
expectation \(d_n\) is finite almost surely, integrable and
\(\mathcal F_{n-1}\)-measurable. Choose a finite measurable version.
Every finite sum \(D_n\) is integrable and
\(\mathcal F_{n-1}\)-measurable, since earlier information fields
are included in that field. The process \(M=Y-D\) is adapted
and integrable, and
\[
 E(M_n\mid\mathcal F_{n-1})
       =E(Y_n\mid\mathcal F_{n-1})-D_n
       =Y_{n-1}+d_n-D_n=M_{n-1}.
\]
Thus it is a martingale.

For any other decomposition \(Y_n=M'_n+D'_n\) of the stated type,
\(D'_n=Y_n-M'_n\) is integrable. Predictability and the martingale
identity imply
\[
           D'_n-D'_{n-1}
               =E(Y_n\mid\mathcal F_{n-1})-Y_{n-1}=d_n
                                    \quad\text{almost surely}.
\]
Starting from \(D'_0=0\), induction gives \(D'_n=D_n\) and
\(M'_n=M_n\). Countably many timewise equalities may be taken
together outside one null set.
The submartingale condition is exactly \(d_n\geq0\) almost surely.
For that case replace any exceptional negative or nonfinite
values separately in each \(\mathcal F_{n-1}\)-measurable version
by zero. The resulting nonnegative finite increments make \(D\)
nondecreasing everywhere while preserving all conditional identities.
Conversely an increasing predictable drift gives
\(E(Y_n\mid\mathcal F_{n-1})=Y_{n-1}+D_n-D_{n-1}\geq Y_{n-1}\).
\end{proof}
\end{annexitem}

\begin{annexitem}{annex:p15.bounded-increment-dichotomy}{thm:p15-bounded-increment-dichotomy}{Lower and upper stopping barriers}
\begin{proof}
Subtract the integrable \(\mathcal F_0\)-measurable \(X_0\).
The process \(Y_n=X_n-X_0\) is a martingale starting at zero
with the same increments. Adding back a finite pathwise constant
changes neither a finite-limit property nor either infinite limit.

Fix a positive integer \(r\) and put
\(T_r=\inf\{n\geq0:Y_n\leq-r\}\).
Before this time \(Y_n>-r\); at the first crossing the fixed
increment bound gives \(Y_{T_r}\geq-r-K\).
Hence \(Y_{T_r\wedge n}+r+K\) is a nonnegative martingale.
It has a finite almost-sure limit.
On \(\{T_r=\infty\}\), therefore, \(Y_n\) converges to a finite limit.
These implications hold together for all positive integer \(r\).

If \(\liminf_nY_n>-\infty\), the path has a finite lower bound:
its tail is bounded below and its finite initial segment consists
of finite numbers. Some integer \(r\) then has \(T_r=\infty\).
Thus \(Y_n\) converges finitely almost surely on this event.
Apply the same argument to \(-Y_n\) to get finite convergence
on \(\{\limsup_nY_n<\infty\}\).
Outside the union of these two events, the lower limit must be
\(-\infty\) and the upper limit \(+\infty\).
This proves the dichotomy.
\end{proof}
\end{annexitem}

\begin{annexitem}{annex:p15.conditional-borel-cantelli-source}{thm:p15-conditional-borel-cantelli-source}{The complete source count dichotomy}
\begin{proof}
Conditional probability gives
\[
 E(\mathbf1_{B_n}-p_n\mid\mathcal F_{n-1})=0,\qquad
                    |\mathbf1_{B_n}-p_n|\leq1.
\]
Thus \(M_n=C_n-P_n\) is an integrable martingale starting at zero
with uniformly bounded increments.
Both \(C_n\) and \(P_n\) are nonnegative and nondecreasing.

On the event that \(M_n\) has a finite limit, their difference
is bounded, so one tends to infinity exactly when the other does.
On the event of two-sided unbounded oscillation,
\(M_n=C_n-P_n\leq C_n\) and \(-M_n=P_n-C_n\leq P_n\)
force both increasing sums to tend to infinity.
The preceding dichotomy exhausts the outcomes up to a null set.
Finally \(C_n\to\infty\) is exactly the event that infinitely
many \(B_n\) occur. This proves the asserted event equality.
No step required the initial sigma-algebra to be trivial.
\end{proof}
\end{annexitem}

\begin{annexitem}{annex:p15.summable-adapted-errors}{thm:p15-summable-adapted-errors}{Stopping before an excessive allowance}
\begin{proof}
Fix a positive integer \(r\), and define
\[
                     T_r=\inf\{n\geq0:\sum_{j=0}^nV_j>r\}.
\]
This is a stopping time, because each partial allowance sum is
known at its indexed time. Put
\[
 C_n^V=\sum_{j=0}^{n-1}\mathbf1_{\{T_r>j\}}V_j,\qquad
 C_n^U=\sum_{j=0}^{n-1}\mathbf1_{\{T_r>j\}}U_j,\qquad
 W_n=Z_{T_r\wedge n}+r-C_n^V .
\]
Empty sums are zero. The increment at the first excessive
allowance is excluded, so \(0\leq C_n^V\leq r\).
Consequently \(W_n\geq0\) and is integrable at every finite time.
On \(\{T_r\leq n\}\) its stopped \(Z\) value is already known;
on \(\{T_r>n\}\) apply the stated conditional inequality.
Subtracting the next allowance gives
\[
                  E(W_{n+1}\mid\mathcal F_n)
                          \leq W_n-\mathbf1_{\{T_r>n\}}U_n.
\]
Thus \(W\) is a nonnegative supermartingale and converges
almost surely to a finite limit.
Also \(C_n^V\) increases to a value at most \(r\), so
\(Z_{T_r\wedge n}=W_n-r+C_n^V\) has a finite limit.

Taking expectations in the last inequality and summing from
zero to \(n-1\) gives
\[
                            EC_n^U\leq EW_0-EW_n\leq EZ_0+r .
\]
Monotone convergence implies \(EC_\infty^U\leq EZ_0+r<\infty\);
hence \(C_\infty^U<\infty\) almost surely.
On \(\{\sum_{j\geq0}V_j\leq r\}\), \(T_r=\infty\), so these
conclusions apply to the original \(Z_n\) and the full sum of \(U_j\).
The countable union over positive integers \(r\) covers
\(\{\sum_jV_j<\infty\}\). This proves both conclusions without
an integrability assumption on the total allowance.
\end{proof}
\end{annexitem}

\begin{annexitem}{annex:p15.predictable-variation-identities}{thm:p15-predictable-variation-identities}{Both compensators and the exact square expansion}
\begin{proof}
Square integrability gives \(E(\Delta X_n)^2<\infty\).
Its conditional expectation has a finite nonnegative
\(\mathcal F_{n-1}\)-measurable version, so every finite
\(A_n=\sum_{j\leq n}a_j\) is predictable, nonnegative and integrable.
The martingale conditional-variance identity gives
\[
                     E(X_n^2\mid\mathcal F_{n-1})
                                           =X_{n-1}^2+a_n.
\]
Therefore \(X_n^2-A_n\) is an integrable martingale.
Likewise the increment of \([X]_n-A_n\) is
\((\Delta X_n)^2-a_n\), with conditional mean zero.
Both start at zero, giving \(EA_n=EX_n^2=E([X]_n)\).

The algebraic identity
\[
           X_j^2-X_{j-1}^2
                     =2X_{j-1}\Delta X_j+(\Delta X_j)^2
\]
sums to the stated discrete square formula.
Cauchy--Schwarz makes each \(X_{j-1}\Delta X_j\) integrable.
The conditional pullout theorem then gives
\[
 E(X_{j-1}\Delta X_j\mid\mathcal F_{j-1})
                              =X_{j-1}E(\Delta X_j\mid\mathcal F_{j-1})=0.
\]
Thus the stochastic sum is an integrable martingale by direct
conditional calculation. This argument supplies the needed
integrability for its unbounded predictable coefficients.
\end{proof}
\end{annexitem}

\begin{annexitem}{annex:p15.variation-stopping-transform}{thm:p15-variation-stopping-transform}{Predictable indicators and transformed second moments}
\begin{proof}
For any stopping time,
\[
             X_{T\wedge n}-X_{T\wedge(n-1)}
                               =\mathbf1_{\{T\geq n\}}\Delta X_n.
\]
The indicator is \(\mathcal F_{n-1}\)-measurable.
The stopped finite-time values are square integrable, since their
squares are at most \(\sum_{j=0}^nX_j^2\).
Conditional expectation of the squared increment therefore equals
\(\mathbf1_{\{T\geq n\}}a_n\).
Summing gives
\[
                  \langle X^T\rangle_n
                    =\sum_{j=1}^n\mathbf1_{\{T\geq j\}}a_j=A_{T\wedge n}.
\]
The analogous pathwise sum of actual squared increments is
\([X]_{T\wedge n}\).

For the bounded-at-each-time predictable transform,
each \(H_n\Delta X_n\) is in \(L^2\); its finite sums form a
square-integrable martingale. Taking its conditional second
increment moment gives
\[
 E(H_n^2(\Delta X_n)^2\mid\mathcal F_{n-1})=H_n^2a_n,
\]
proving the transform formula.
For each finite \(n\), expectation of the stopped compensated
variation is zero, so
\(E([X]_{T\wedge n})=EA_{T\wedge n}\).
Both sides are expectations of increasing nonnegative sequences.
Monotone convergence gives \(E([X]_T)=EA_T\), allowing infinity.
No passage of \(X_{T\wedge n}^2\) to an unbounded-time square
was used in this last argument.

The absence of an unbounded-time terminal-square conclusion has
a concrete example independent of the later square-root estimate.
Let \(S\) be a fair-sign walk starting at zero and \(H\) its first
hit of \(-1\), finite almost surely by the earlier hitting theorem.
Its variance process is \(A_n=n\), whereas \(S_H^2=1\).
If \(EH<\infty\), then \(|S_{H\wedge n}|\leq H\) supplies an
integrable bound. Bounded optional sampling and dominated convergence
would give \(ES_H=0\), contradicting \(S_H=-1\).
Thus \(EA_H=EH=\infty\) while \(ES_H^2=1\).
\end{proof}
\end{annexitem}

\begin{annexitem}{annex:p15.quadratic-maximal-bound}{thm:p15-quadratic-maximal-bound}{The maximal square estimate and its finite-expectation consequence}
\begin{proof}
At finite horizon the \(L^2\) maximal inequality gives
\[
                    E\max_{0\leq j\leq n}|X_j|^2
                                 \leq4EX_n^2=4EA_n.
\]
Both nonnegative sides have increasing limits; monotone
convergence proves \(E(X^*)^2\leq4EA_\infty\).
If \(EA_\infty<\infty\), then
\(\sup_nEX_n^2=\sup_nEA_n\leq EA_\infty\).
The proved \(L^2\) convergence theorem yields an almost-sure
and mean-square limit.
Continuity of the \(L^2\) norm under that convergence gives
\(EX_\infty^2=\lim_nEX_n^2=\lim_nEA_n=EA_\infty\).
\end{proof}
\end{annexitem}

\begin{annexitem}{annex:p15.finite-variation-convergence}{thm:p15-finite-variation-convergence}{Localization by the next known variance value}
\begin{proof}
Because \(A_{n+1}\) is \(\mathcal F_n\)-measurable,
\[
                 \{T_c\leq n\}
                    =\bigcup_{j=0}^n\{A_{j+1}>c\}\in\mathcal F_n.
\]
If \(T_c=k>0\), its minimality gives \(A_k\leq c\);
if \(T_c=0\), \(A_0=0\).
Before the stopping time, all the corresponding \(A_n\) values
are at most \(c\). If the time is infinite, this is true
for every \(n\). Thus \(A_{T_c\wedge n}\leq c\) in every case.
The stopped variation identity and expectation give
\[
                             EX_{T_c\wedge n}^2
                                  =EA_{T_c\wedge n}\leq c.
\]
Hence the stopped martingale converges almost surely and in
\(L^2\), with a finite limit.
For positive integer \(c\), these convergence assertions hold
together outside one null set.
On \(\{A_\infty<\infty\}\), choose an integer \(c\geq A_\infty\).
Then \(T_c=\infty\), so the stopped sequence is the original
sequence. This proves the eventwise finite-limit assertion.
It does not assert global integrability of that limit merely
from pathwise finiteness of \(A_\infty\).
\end{proof}
\end{annexitem}

\begin{annexitem}{annex:p15.bounded-jump-variation-dichotomy}{thm:p15-bounded-jump-variation-dichotomy}{The bounded-overshoot converse}
\begin{proof}
The forward implication from finite \(A_\infty\) to finite
convergence was just proved.
For the converse fix a positive integer \(r\) and let
\(T_r=\inf\{n\geq0:|X_n|>r\}\).
The fixed increment bound gives
\(|X_{T_r\wedge n}|\leq r+K\) almost surely for every \(n\):
at a crossing use the preceding value and one increment,
and otherwise use the unexited bound.
Consequently
\[
                   EA_{T_r\wedge n}=EX_{T_r\wedge n}^2
                                              \leq(r+K)^2.
\]
Monotone convergence yields \(EA_{T_r}\leq(r+K)^2\), so
\(A_{T_r}<\infty\) almost surely.
These conclusions hold simultaneously for all integer \(r\).
Every real sequence with a finite limit is bounded, including
its finite initial segment. On a convergent path some \(r\)
therefore has \(T_r=\infty\), giving \(A_\infty=A_{T_r}<\infty\).
This proves the equality of the two events up to null sets.
The bounded-increment convergence-or-oscillation theorem then
gives the stated two infinite limits on the complementary event.
\end{proof}
\end{annexitem}

\begin{annexitem}{annex:p15.variance-clock-normalization}{thm:p15-variance-clock-normalization}{The right-endpoint variance sum and pathwise Kronecker}
\begin{proof}
Since \(f\geq1\), the coefficients \(H_j=1/f(A_j)\) are bounded
by one. They are predictable because \(A_j\) is known in
\(\mathcal F_{j-1}\). The transformed process \(Y\) is a
square-integrable martingale, with variance process
\[
                           B_n=\sum_{j=1}^n
                                    \frac{A_j-A_{j-1}}{f(A_j)^2}.
\]
For every path, monotonicity of \(f\) gives
\[
            \frac{A_j-A_{j-1}}{f(A_j)^2}
                         \leq\int_{[A_{j-1},A_j)}f(t)^{-2}\,dt.
\]
If the endpoints coincide, both sides are zero.
Otherwise the inequality follows from \(f(t)\leq f(A_j)\)
on the integration interval; continuity of \(f\) is unnecessary.
The intervals are disjoint and exhaust \([0,A_\infty)\).
Summing their nonnegative integrals gives
\(B_\infty\leq I_f<\infty\).
The finite-variation theorem now makes \(Y_n\) converge almost
surely. In fact \(EB_\infty\leq I_f\) also gives \(L^2\) convergence.

A finite nondecreasing \(f\) with integrable reciprocal square
must tend to infinity. If it had a finite upper limit \(L\),
then \(f(t)^{-2}\geq L^{-2}>0\) for all \(t\), contradicting
the finite integral.
Thus on \(\{A_\infty=\infty\}\), the pathwise weights
\(b_j=f(A_j)\) are positive, nondecreasing and divergent.
On the common event where the transformed series converges,
apply Kronecker's lemma to \(x_j=\Delta X_j\) and these
numerical weights. It gives
\[
                    \frac1{f(A_n)}\sum_{j=1}^n\Delta X_j\longrightarrow0.
\]
The sum is \(X_n\) because \(X_0=0\), proving the result.
\end{proof}
\end{annexitem}

\begin{annexitem}{annex:p15.variance-logarithmic-rate}{thm:p15-variance-logarithmic-rate}{An everywhere defined version of the source scale}
\begin{proof}
The constant branch is \(\sqrt e\geq1\); the positive tail
\(t(\log t)^{1+\epsilon}\) is increasing for \(t\geq e\)
and equals \(e\) at \(e\). Thus the two branches join
continuously and give a finite nondecreasing Borel function.
Their reciprocal-square integrals are
\[
 \int_0^e e^{-1}\,dt=1,\qquad
 \int_e^\infty\frac{dt}{t(\log t)^{1+\epsilon}}
        =\int_1^\infty u^{-1-\epsilon}\,du=\epsilon^{-1}.
\]
The normalization theorem applies.
On the event \(A_\infty=\infty\), eventually \(A_n>e\),
where the scale is precisely the displayed source tail expression.

For independent centered increments, conditional independence gives
\(E((\Delta X_j)^2\mid\mathcal F_{j-1})=\sigma_j^2\).
Hence \(A_n=\sum_{j\leq n}\sigma_j^2\), proving the independent-sum
specialization. If the variances are identically \(\sigma^2>0\),
\(A_n=n\sigma^2\), and the logarithmic scale is asymptotic to
\(\sigma\sqrt{n(\log n)^{1+\epsilon}}\); this follows from
\(\log(n\sigma^2)/\log n\to1\).
The zero-variance case has no divergent variance event and makes
all centered increments zero almost surely.
\end{proof}
\end{annexitem}

\begin{annexitem}{annex:p15.conditional-borel-cantelli-ratio}{thm:p15-conditional-borel-cantelli-ratio}{Both variance regimes of the conditional count law}
\begin{proof}
The finite-time compensated count is bounded in absolute value by \(n\),
so it is square integrable. Conditional expansion of its squared
increment gives
\[
 \begin{aligned}
 E((\mathbf1_{B_n}-p_n)^2\mid\mathcal F_{n-1})
       &=p_n-2p_n^2+p_n^2=p_n(1-p_n).
 \end{aligned}
\]
Write \(V_n=\sum_{j\leq n}p_j(1-p_j)\) for this variance process.
Then \(0\leq V_n\leq P_n\).

On \(\{V_\infty<\infty\}\), the finite-variation theorem makes
\(M_n\) converge to a finite limit. If also \(P_\infty=\infty\),
this gives \(M_n/P_n\to0\).
On \(\{V_\infty=\infty\}\), also \(P_\infty=\infty\).
The function \(f(t)=\max(1,t)\) has reciprocal-square integral
\(1+1=2\). The normalization theorem gives
\(M_n/f(V_n)\to0\) almost surely on this event.
Eventually \(V_n>1\), and
\[
                        \frac{|M_n|}{P_n}\leq\frac{|M_n|}{V_n}\longrightarrow0.
\]
These two cases exhaust \(\{P_\infty=\infty\}\).
Since \(C_n/P_n=1+M_n/P_n\) whenever \(P_n>0\), the ratio follows.

On \(\{P_\infty<\infty\}\), \(V_\infty<\infty\), so
\(C_n=M_n+P_n\) has a finite limit. It is an increasing
integer-valued count and therefore is eventually constant.
Thus only finitely many events occur there.
This supplies a second proof of the finite-event part while
retaining the source's separate convergence-or-oscillation route above.
\end{proof}
\end{annexitem}

\begin{annexitem}{annex:p15.friedman-urn-limit}{thm:p15-friedman-urn-limit}{All source cases and the zero same-colour boundary}
\begin{proof}
Use independent uniform inputs to realize each draw with its stated
conditional probability, and let \(\mathcal F_n\) contain their
first \(n\) values. Put \(c=a+b>0\), \(t_0=g+r\).
The displayed colour counts follow by counting each type of draw.
They remain positive and have deterministic total \(t_n=t_0+cn\).
In particular
\[
                           g_n\geq\frac g{t_0+cn},\qquad
                           r_n\geq\frac r{t_0+cn}.
\]
Both sums of successive draw probabilities diverge.
The conditional count law therefore gives, almost surely,
\[
 \frac{D_n}{\sum_{j=1}^ng_{j-1}}\longrightarrow1,\qquad
 \frac{n-D_n}{\sum_{j=1}^nr_{j-1}}\longrightarrow1.
\]
The denominators are positive and no larger than \(n\).
Hence the differences between each observed frequency and
its average conditional probability tend to zero.
If a bounded sequence has upper limit at most \(x\), its
Cesaro averages have upper limit at most \(x\): bound its
tail by \(x+\epsilon\), let the finite head vanish after division
by \(n\), then let \(\epsilon\downarrow0\).
Thus any bounds on the upper limits of \(g_n,r_n\) give the
same bounds on the respective draw frequencies.

If \(a=b>0\), both colours receive \(a\) at every draw and
\(g_n=(g+an)/(t_0+2an)\to1/2\).
If \(a>b\), set \(\alpha=(a-b)/(a+b)\in(0,1)\).
The exact count formula, divided by \(n\), shows that an upper
bound \(x\) for \(\limsup g_n\) improves to
\[
                      \limsup g_n\leq\frac{b+(a-b)x}{a+b}.
\]
The same statement holds for the red fraction by interchanging colours.
Beginning with \(x_0=1\), iteration gives
\(x_j=1/2+\alpha^j/2\), so both colour upper limits are at most
\(1/2\). Since they sum to one, both fractions converge to \(1/2\).

If \(0<a<b\), use the number of red draws in the green formula:
\[
                G_n=g+an+(b-a)(n-D_n).
\]
If \(\limsup g_n\leq x\) and \(\limsup r_n\leq y\), the two
improved bounds are
\[
          \limsup g_n\leq\frac{a+(b-a)y}{a+b},\qquad
          \limsup r_n\leq\frac{a+(b-a)x}{a+b}.
\]
Starting with equal bounds one, they remain equal at each
iteration and take the values
\(1/2+\frac12((b-a)/(a+b))^j\).
Their limit is \(1/2\), proving this case.

It remains to prove \(a=0\), where that last multiplier equals one.
Put \(Y_n=g_n-1/2\) and
\(\epsilon_{n+1}=I_{n+1}-g_n\).
The exact update gives
\[
             Y_{n+1}
                =\left(1-\frac{2b}{t_{n+1}}\right)Y_n
                       -\frac b{t_{n+1}}\epsilon_{n+1}.
\]
Conditional on \(\mathcal F_n\), the last variable has mean
zero and variance \(g_n(1-g_n)\leq1/4\).
At all sufficiently large deterministic \(n\), \(t_{n+1}\geq2b\);
therefore, with \(Z_n=Y_n^2\),
\[
 \begin{aligned}
 E(Z_{n+1}\mid\mathcal F_n)
   &\leq\left(1-\frac{4b}{t_{n+1}}
                    +\frac{4b^2}{t_{n+1}^2}\right)Z_n
                        +\frac{b^2}{4t_{n+1}^2}\\
   &\leq Z_n-\frac{2b}{t_{n+1}}Z_n+\frac{b^2}{4t_{n+1}^2}.
 \end{aligned}
\]
Here \(0\leq Z_n\leq1/4\), so all terms are integrable.
The deterministic allowance series \(\sum_n b^2/(4t_{n+1}^2)\)
converges. Apply the summable-error theorem to this tail,
with decrement \(2bZ_n/t_{n+1}\).
It gives a finite limit for \(Z_n\) and
\(\sum_n Z_n/t_{n+1}<\infty\) almost surely.
But \(\sum_n1/t_{n+1}=\infty\).
A positive limiting value of \(Z_n\) would contradict this
finite weighted sum. Hence \(Z_n\to0\), proving \(g_n\to1/2\)
also when \(a=0\).

In every case, Cesaro averaging now gives
\(n^{-1}\sum_{j=1}^ng_{j-1}\to1/2\).
The vanishing frequency difference established at the start
gives \(D_n/n\to1/2\). The two divergent conditional-count
sums also prove infinitely many draws of each colour.
\end{proof}
\end{annexitem}

\begin{annexitem}{annex:p15.square-root-variation-bound}{thm:p15-square-root-variation-bound}{The complete stopped tail integral}
\begin{proof}
Fix \(a>0\) and take \(T=T_{a^2}\), the stop before the next
variance value exceeds \(a^2\).
Its stopped variation satisfies
\(A_{T\wedge n}\leq\min(A_\infty,a^2)\).
Also \(\{T<\infty\}=\{A_\infty>a^2\}\):
a strict exceedance of the increasing limit occurs at some
finite variance value, and conversely any stopped exceedance
makes the limit strictly larger.
On \(\{T=\infty\}\) the stopped and original paths coincide.
Therefore
\[
 P(X^*>a)\leq P(A_\infty>a^2)
                       +P\left(\sup_m|X_{T\wedge m}|>a\right).
\]
At any finite horizon, Doob's probability inequality for the
squared stopped martingale gives
\[
 P\left(\max_{m\leq n}|X_{T\wedge m}|>a\right)
      \leq a^{-2}EX_{T\wedge n}^2
      =a^{-2}EA_{T\wedge n}
      \leq a^{-2}E(A_\infty\wedge a^2).
\]
Let \(n\to\infty\); the strict finite-maximum events increase
to the displayed strict supremum event.
This proves for every \(a>0\)
\[
              P(X^*>a)\leq P(A_\infty>a^2)
                                      +a^{-2}E(A_\infty\wedge a^2).
\]

Integrate over \(a>0\). The first right-hand integral is
\(E\sqrt{A_\infty}\) by the tail identity.
For the second, the full source interchange is justified by
nonnegativity:
\[
 \begin{aligned}
 \int_0^\infty a^{-2}E(A_\infty\wedge a^2)\,da
   &=\int_0^\infty a^{-2}\int_0^{a^2}P(A_\infty>b)\,db\,da\\
   &=\int_0^\infty P(A_\infty>b)
                       \left(\int_{\sqrt b}^\infty a^{-2}\,da\right)db\\
   &=\int_0^\infty b^{-1/2}P(A_\infty>b)\,db
     =2E\sqrt{A_\infty}.
 \end{aligned}
\]
The expressions with \(b^{-1/2}\) concern \(b>0\);
the single endpoint has zero Lebesgue measure.
All interchanges and equalities remain valid as nonnegative
extended integrals, including when \(A_\infty\) can be infinite.
The tail integral of \(X^*\) on the left is \(EX^*\).
Adding the two terms proves the coefficient three.
\end{proof}
\end{annexitem}

\begin{annexitem}{annex:p15.square-root-stopping}{thm:p15-square-root-stopping}{The integrable stopped maximum and conditional sampling}
\begin{proof}
The stopped process has variation \(A_{T\wedge n}\) and limiting
variation \(A_T\).
Apply the square-root maximal theorem to it to obtain the
integrable bound
\[
                      Z=\sup_n|X_{T\wedge n}|,\qquad
                      EZ\leq3E\sqrt{A_T}<\infty.
\]
Since \(T\) is finite almost surely, its stopped sequence is
eventually equal to \(X_T\). It is dominated by \(Z\), so
dominated convergence gives \(L^1\) convergence.
Every bounded stopped value has expectation zero, hence \(EX_T=0\).

Let \(S\leq T\) pointwise and \(A\in\mathcal F_S\).
For an integer \(n\), put \(B_n=A\cap\{S\leq n\}\).
This event is in \(\mathcal F_{S\wedge n}\):
for \(m<n\), intersecting with \(\{S\wedge n\leq m\}\)
gives \(A\cap\{S\leq m\}\in\mathcal F_m\);
for \(m\geq n\), the event is \(B_n\in\mathcal F_n\subseteq\mathcal F_m\).
Bounded optional sampling at \(S\wedge n,T\wedge n\) therefore gives
\[
                  E(X_{T\wedge n};B_n)
                           =E(X_{S\wedge n};B_n)=E(X_S;B_n).
\]
The time \(S\) is finite almost surely and \(|X_S|\leq Z\).
Its stopped-field measurability follows from the event partition
over its finite values, as for earlier stopping results.
Let \(n\to\infty\). The indicators of \(B_n\) tend to that of
\(A\) almost surely, and both products are dominated by \(Z\).
Thus \(E(X_T;A)=E(X_S;A)\) for every \(A\in\mathcal F_S\),
which identifies \(E(X_T\mid\mathcal F_S)=X_S\).

If \(EA_T<\infty\), apply the quadratic maximal and \(L^2\)
convergence theorem to the same stopped process.
Its limit is \(X_T\), and the norm identity gives \(EX_T^2=EA_T\).
For nonzero \(X_0\), the process \(X-X_0\) is a martingale
starting at zero with the same increments and variance process.
Applying the preceding statements and then adding back the
\(\mathcal F_0\)-measurable \(X_0\) proves the stated extension.

For a fair-sign walk the variance process is \(A_n=n\),
so the stopped variance is \(A_T=T\).
The first hit \(H\) of \(-1\) is finite almost surely by
applying the earlier fair-walk hitting argument to \(S_n+1\).
If \(E\sqrt H\) were finite, the criterion just proved would
give \(ES_H=0\). But \(S_H=-1\) almost surely.
This contradiction proves \(E\sqrt H=\infty\).
\end{proof}
\end{annexitem}

\begin{annexitem}{annex:p15.variance-clock-examples}{thm:p15-variance-clock-examples}{Every variance and integrability claim in the examples}
\begin{proof}
For the first sequence, each increment has mean zero and
second moment \(n^2(1/n^2)=1\).
Independence gives conditional second moment one, hence \(A_n=n\).
The probabilities of nonzero increments sum to
\(\sum_n1/n^2<\infty\).
Borel--Cantelli gives only finitely many nonzero increments
almost surely. Their finite sum then becomes constant, and the
sum of their actual squares is finite on that same event.
All finite-time second moments are finite.
Thus neither a finite limit nor finite realized total variation
implies finite predictable total variation without further hypotheses.

The positive normalizing series for \(c\) in the second example
converges by the power-series integral comparison, so it defines
a probability law on the positive integers.
The same comparison gives
\[
                 EN=c\sum_{k\geq1}k^{-3/4}=\infty,\qquad
                 E\sqrt N=c\sum_{k\geq1}k^{-5/4}<\infty.
\]
Because \(N\) is independent of the walk, revealing it in
\(\mathcal F_0\) preserves the martingale property of the walk.
It is a stopping time in the enlarged filtration.
The stopped increments are \(\mathbf1_{\{N\geq n\}}\xi_n\),
so their conditional second moments are
\(\mathbf1_{\{N\geq n\}}\). Thus \(A_n=N\wedge n\) and
\(A_\infty=N\).
The square-root stopping criterion gives a finite first absolute
moment and expectation zero for \(S_N\).
Independence and Tonelli for the nonnegative square give
\[
                         ES_N^2=\sum_{k\geq1}P(N=k)ES_k^2
                                             =\sum_{k\geq1}kP(N=k)=\infty.
\]
This verifies the strict difference between the two moment conditions.

In the third example the only nonzero conditional increment
variance is \(a_1=1\). Hence \(A_n=1\) from time one onward.
The value remains the first sign, and division by the finite
positive \(f(1)\) cannot make its limit zero.
\end{proof}
\end{annexitem}

\begin{annexitem}{annex:p15.randomly-bounded-paths}{thm:p15-randomly-bounded-paths}{Rare compensation, bounded paths and no probability limit}
\begin{proof}
Let \(\mathcal F_n=\sigma(\eta_j,\zeta_j:1\leq j\leq n)\),
with trivial \(\mathcal F_0\).
The recursive values are finite and adapted.
Induction gives integrability, because
\[
                 E|X_n|=(1-\epsilon_n)E|\zeta_n|+E|X_{n-1}|
                               =2p(1-\epsilon_n)+E|X_{n-1}|.
\]
The formula follows by the disjoint cases \(\eta_n=0,1\)
and independence of the new inputs from the past.
Summing it yields
\(E|X_n|=2p(n-1+2^{-n})\).
The conditional mean is
\[
 E(X_n\mid\mathcal F_{n-1})
        =(1-\epsilon_n)E\zeta_n
                       +\epsilon_nX_{n-1}/\epsilon_n=X_{n-1}.
\]
The analogous square recursion is
\[
                        EX_n^2=2p(1-\epsilon_n)
                                             +\epsilon_n^{-1}EX_{n-1}^2.
\]
It proves square integrability at every finite time.

Since \(\sum_n\epsilon_n<\infty\), only finitely many
\(\eta_n=1\) occur almost surely.
After their last occurrence the recursion gives \(X_n=\zeta_n\).
The finite earlier values are finite numbers, so the full path
is bounded. Each of the three independent-input values has
positive probability, and Borel--Cantelli gives infinitely
many occurrences of each. These properties persist after
any finite initial segment, proving all three almost-sure
infinite-occurrence assertions and failure of almost-sure convergence.

Under the displayed coupling, \(P(X_n\ne\zeta_n)\leq\epsilon_n\).
It follows immediately that the probabilities at \(-1,0,1\)
have the stated limits.
For probability convergence, put
\[
                         q=1-[(1-2p)^2+2p^2]=4p-6p^2>0.
\]
Independent successive \(\zeta\)'s differ with probability \(q\),
and every nonzero difference has absolute value at least one.
Except on the union of the two rare compensation events,
the successive \(X\)'s equal these inputs. Hence
\[
                  P(|X_{n+1}-X_n|>1/2)
                                   \geq q-\epsilon_n-\epsilon_{n+1}.
\]
The lower bound tends to \(q>0\).
Convergence in probability would make these successive
differences tend to zero in probability by the triangle inequality,
so no such limit exists.

Finally, for any deterministic proposed increment bound \(K\),
choose \(n\geq2\) with \(2^n-1>K\).
The event \(\{\eta_{n-1}=0,\zeta_{n-1}=1,\eta_n=1\}\)
has probability \((1-\epsilon_{n-1})p\epsilon_n>0\).
There \(X_{n-1}=1\) and \(X_n=2^n\), violating that bound.
This proves the exact distinction from uniformly bounded increments.
\end{proof}
\end{annexitem}

\begin{annexitem}{annex:p15.urn-boundary-calculations}{thm:p15-urn-boundary-calculations}{The same-colour limit and the random initial environment}
\begin{proof}
When \(a=b>0\), every draw adds \(a\) to each colour,
giving the stated deterministic fraction directly.

For the same-colour urn with \(a=1,b=0,g=r=1\), write \(G,R\)
for the current counts and \(t=G+R\). The expected next fraction is
\[
            \frac Gt\frac{G+1}{t+1}
                   +\frac Rt\frac G{t+1}=\frac Gt.
\]
Thus its fraction is a bounded martingale and has an almost-sure
limit. A specified word of \(n\) draws containing \(m\) greens
and \(n-m\) reds has probability
\[
                              \frac{m!(n-m)!}{(n+1)!}.
\]
The denominator is the product \(2\cdot3\cdots(n+1)\);
the green and red numerators separately run through
\(1,\ldots,m\) and \(1,\ldots,n-m\), regardless of their interleaving.
There are \(\binom nm\) such words, so
\(P(G_n=m+1)=1/(n+1)\), \(0\leq m\leq n\).
Its fractions therefore have the uniform law on
\(\{1/(n+2),\ldots,(n+1)/(n+2)\}\).
For every continuous \(h\) on \([0,1]\), the corresponding
finite average tends to \(\int_0^1h(t)\,dt\) by uniform
continuity and the equal-mesh Riemann sums.
Bounded convergence also makes that average tend to the
expectation of \(h\) of the almost-sure limit.
Uniqueness from continuous tests identifies the limit as
uniform on \((0,1)\). In particular it is not almost surely \(1/2\).

For \(B_n=\Omega\), every \(p_n=1\), so \(C_n=P_n=n\),
\(M_n=0\) and its variance is zero.
For the random-environment construction use
\(\mathcal F_n=\sigma(\Theta,U_1,\ldots,U_n)\).
Independence of the next uniform input and its uniform CDF
give \(p_n=\Theta/(n+1)\). Thus
\[
                            P_n=\Theta\sum_{j=1}^n\frac1{j+1}.
\]
On \(\{\Theta>0\}\) this diverges, and the normalized
conditional-count theorem gives \(C_n/P_n\to1\).
The harmonic sum divided by \(\log n\) tends to one by
the elementary integral bounds for \(\sum 1/j\).
Hence \(C_n/\log n\to\Theta\) on this event.
On \(\{\Theta=0\}\), each \(U_n>0\), so no event occurs
and the conclusion is zero directly. This uses \(n\geq2\)
and never divides by \(\Theta\) on its zero event.
\end{proof}
\end{annexitem}

\begin{annexitem}{annex:p15.decomposition-rate-boundaries}{thm:p15-decomposition-rate-boundaries}{Three explicit missing-hypothesis calculations}
\begin{proof}
For the Bernoulli example let \(\mathcal F_0\) be trivial and
\(\mathcal F_n=\sigma(B)\) for \(n\geq1\).
The conditional drift at the first step is \(EB=1/2\);
later drifts vanish. The stated \(D\) is predictable and
increasing, and \(M_1=B-1/2\), held constant afterward,
is a martingale starting at zero.
The alternative second part \(Y_n\) is also increasing and
adapted, but \(Y_1=B\) is not measurable in the trivial
\(\mathcal F_0\). Both decompositions have zero initial
second part, establishing the necessity of predictability for uniqueness.

For \(X_n=c\ne0\), every increment and hence every \(a_n\)
is zero, whereas \(X^*=|c|>0\).
This gives both failures without the initial-value condition.

For the fair-sign walk, independence and symmetry give
\[
                         ES_n^2=n,\qquad
                         ES_n^4=n+6\binom n2=3n^2-2n.
\]
In the fourth-power expansion, all terms with an odd power
of some independent sign have zero expectation; the remaining
terms are the \(n\) fourth powers and the six arrangements
of each pair of squared signs.
Put \(Z=S_n^2\). Splitting at \(n/2\) and applying
Cauchy--Schwarz gives
\[
 n=EZ\leq n/2+
             (EZ^2)^{1/2}P(Z>n/2)^{1/2}.
\]
Therefore
\[
                 P(Z>n/2)\geq
                     \frac{n^2}{4(3n^2-2n)}\geq\frac1{12}.
\]
This is the asserted strict-threshold bound.
Since \(A_n=n\), the proposed scale \(f(A_n)\) is \(\sqrt n\)
for \(n\geq1\); the lower bound prevents convergence to zero
in probability. Finally its reciprocal-square integral
contains \(\int_1^\infty t^{-1}\,dt=\infty\),
so precisely the required integrability hypothesis is absent.
\end{proof}
\end{annexitem}
